\documentclass[11pt,a4paper]{article}
\usepackage[english]{babel}
\usepackage[utf8]{inputenc}
\usepackage[T1]{fontenc}
\usepackage{amsmath,amssymb,bm}
\usepackage{geometry}
\usepackage{xcolor}
\usepackage{booktabs}
\usepackage{hyperref}
\geometry{margin=2.5cm}
\hypersetup{colorlinks=true,linkcolor=blue,urlcolor=blue}

\newcommand{\BHF}{B_{\mathrm{HF}}}
\newcommand{\DEQ}{\Delta E_Q}
\newcommand{\dd}{\mathrm{d}}

\title{Two-dimensional distributions $P(\BHF,\DEQ)$ in M\"{o}ssbauer spectra}
\author{Fitbauer -- mathematical implementation note}
\date{\today}

\begin{document}
\maketitle

\section{Motivation}

One-dimensional distributions $P(\BHF)$, $P(\DEQ)$, and $P(\delta)$ are
adequate when the dominant source of disorder is, respectively, magnetic,
quadrupolar, or isomer-shift. In amorphous materials, disordered ferrites,
glasses, nanoparticles, or phases with multiple local environments, two
parameters may vary simultaneously. In that case a two-dimensional distribution
can be formulated:
\begin{equation}
  P(x,y)\ge 0,\qquad x,y\in\{\BHF,\DEQ,\delta\}.
\end{equation}
The cases implemented in the GUI are $P(\BHF,\DEQ)$, $P(\delta,\DEQ)$, and
$P(\BHF,\delta)$. The main interest is to study not only the marginals but also
the possible correlation between the two distributed quantities. The case
$P(\delta,\DEQ)$ is especially relevant for broad paramagnetic doublets in
glasses, silicates, and amorphous phases.

\section{Spectral model}

Let $v_k$ be the velocity channel and $S(v;B,Q,\delta,\Gamma)$ the unit-depth
absorption sextet. On a grid $x_i$, $y_j$ the model linear in the weights is
\begin{equation}
  T_k = b_0 + b_1 v_k - \sum_{i=1}^{N_x}\sum_{j=1}^{N_y}
  P_{ij}\,S(v_k;B_{ij},Q_{ij},\delta_{ij},\Gamma).
\end{equation}
Non-distributed parameters are kept global. For example, in $P(\delta,\DEQ)$
a global $B$ is fixed or fitted; if $B=0$ the pattern represents a set of
magnetically collapsed doublets.
In matrix form,
\begin{equation}
  \bm y \simeq X\bm z,
\end{equation}
where $\bm z$ contains background, slope, and the flattened weights $P_{ij}$.
The columns associated with the distribution are $-S(v_k;B_i,Q_j)$ because
absorption reduces transmission.

\section{2D regularisation}

The inverse problem is strongly under-determined. Therefore the curvature of
the distribution is penalised in both directions:
\begin{equation}
  \Phi(\bm z)=\|W(\bm y-X\bm z)\|^2
  +\alpha_B\|D_B^2 P\|^2
  +\alpha_Q\|D_Q^2 P\|^2,
\end{equation}
subject to
\begin{equation}
  P_{ij}\ge 0.
\end{equation}
Here $W$ is the statistical weight of the channels, $D_B^2$ applies second
differences along $\BHF$, and $D_Q^2$ applies second differences along
$\DEQ$. In the initial implementation a rectangular grid is used and the
operators are written via Kronecker products:
\begin{align}
  L_B &= D_B^2\otimes I_Q,\\
  L_Q &= I_B\otimes D_Q^2.
\end{align}
The augmented problem solved by bounded least squares is
\begin{equation}
  \begin{bmatrix}
    WX\\
    \sqrt{\alpha_B}\,L_B\\
    \sqrt{\alpha_Q}\,L_Q
  \end{bmatrix}\bm z
  \simeq
  \begin{bmatrix}
    W\bm y\\
    \bm 0\\
    \bm 0
  \end{bmatrix},
  \qquad P_{ij}\ge 0.
\end{equation}

\section{Marginals, moments, and correlation}

Once $P_{ij}$ has been fitted, the marginal distributions are
\begin{align}
  P_B(B_i) &= \int P(B_i,Q)\,\dd Q \simeq \sum_j P_{ij}\,\Delta Q_j,\\
  P_Q(Q_j) &= \int P(B,Q_j)\,\dd B \simeq \sum_i P_{ij}\,\Delta B_i.
\end{align}
On a non-uniform grid appropriate integration weights must be used. The
normalised two-dimensional probability is defined by
\begin{equation}
  p_{ij}=\frac{P_{ij}}{\int\!\int P(B,Q)\,\dd B\,\dd Q}.
\end{equation}
With $p_{ij}$ the diagnostic moments are computed:
\begin{align}
  \langle B\rangle &= \sum_{ij} p_{ij} B_i\,\Delta B_i\Delta Q_j, &
  \langle Q\rangle &= \sum_{ij} p_{ij} Q_j\,\Delta B_i\Delta Q_j,\\
  \sigma_B^2 &= \sum_{ij} p_{ij}(B_i-\langle B\rangle)^2\Delta B_i\Delta Q_j, &
  \sigma_Q^2 &= \sum_{ij} p_{ij}(Q_j-\langle Q\rangle)^2\Delta B_i\Delta Q_j,
\end{align}
and the apparent correlation
\begin{equation}
  \rho_{BQ}=\frac{\sum_{ij}p_{ij}(B_i-\langle B\rangle)(Q_j-\langle Q\rangle)
  \Delta B_i\Delta Q_j}{\sigma_B\sigma_Q}.
\end{equation}
A tilted cloud in the $B$--$Q$ map may suggest correlation between field and
quadrupole splitting, but that interpretation is only valid if the solution is
stable.

\section{Simultaneous sharp components}

The implementation allows sharp components to be added to the 2D problem. The
model extends to
\begin{equation}
  T_k=b_0+b_1v_k-\sum_{ij}P_{ij}S_{ij}(v_k)-\sum_m a_m C_m(v_k)-F_k,
\end{equation}
where $C_m$ are discrete components with free amplitude $a_m\ge0$ and $F_k$ is
the known absorption of components with fixed depth. The sharp columns are not
regularised; they compete with the distribution only through the residual. This
is useful for separating a known discrete phase from a disordered background,
but it can also transfer weight artificially between the map and the sharp
components if they overlap.

\section{Regularisation L-surface in 2D}

In 1D an L-curve is used. In 2D there are two regularisations, so the analogue
is a surface in $(\alpha_B,\alpha_Q)$. For each pair the following are computed:
\begin{align}
  R(\alpha_B,\alpha_Q)&=\|W(\bm y-X\bm z)\|,\\
  G(\alpha_B,\alpha_Q)&=\sqrt{\|D_B^2P\|^2+\|D_Q^2P\|^2},\\
  \mathrm{RMS}(\alpha_B,\alpha_Q)&=\sqrt{\langle r_k^2\rangle}.
\end{align}
The GUI shows a map of RMS versus $\log_{10}\alpha_B$ and $\log_{10}\alpha_Q$
and a projection of $R$ versus $G$. Two points are proposed: the minimum of
RMS and a geometric trade-off between residual and roughness. Neither is an
automatic answer; the stability of the map must be verified visually.

\section{Risk of overfitting}

\textbf{Essential warning:} a 2D distribution can apparently fit very well
and carry no physical meaning. For example, a $30\times30$ grid contains 900
weights, often more than the effective number of independent points in the
spectrum. With weak regularisation, the map can absorb:
\begin{itemize}
  \item statistical noise,
  \item folding or velocity-scale errors,
  \item poorly chosen linewidths,
  \item instrumental background,
  \item omitted discrete components,
  \item dynamic relaxation that should be modelled differently.
\end{itemize}
Therefore, a low residual or a visually attractive map do not by themselves
prove the existence of a physical distribution $P(\BHF,\DEQ)$.

\section{Best practices}

Before interpreting the map it is recommended to:
\begin{enumerate}
  \item compare with simpler models: discrete sextets, $P(\BHF)$,
  $P(\DEQ)$, and sharp components;
  \item vary $\alpha_B$ and $\alpha_Q$ over several orders of magnitude;
  \item repeat with a different number of bins;
  \item verify that the marginals and main features are stable;
  \item fix $\delta$ and $\Gamma$ if known independently;
  \item check that the residual retains no systematic structure;
  \item compare with external information: TEM, magnetometry, structure,
  chemistry, or literature.
\end{enumerate}
For publication, the 2D model should be justified by showing that a 1D or
discrete model does not explain the data satisfactorily and that the 2D
distribution is robust against reasonable changes in regularisation.

\section{Implementation status}

The Fitbauer implementation includes the backend:
\begin{center}
\texttt{mossbauer\_distribution.fit\_bhf\_quad\_distribution()}
\end{center}
and integration in the GUI as modes \textbf{P(BHF,$\DEQ$) 2D},
\textbf{P($\delta$,$\DEQ$) 2D}, \textbf{P(BHF,$\delta$) 2D}, and as the
1D distribution \textbf{P($\delta$)}. The backend returns the centres of both
axes, the weight matrix, the normalised probability, the fitted curve, the
residual, the background, the slope, the values of $\alpha_x$, $\alpha_y$,
marginals, moments, apparent correlation, and identifiability warnings. The GUI
shows a heatmap with marginals, allows exporting the map to TSV, includes an
$\alpha_x$--$\alpha_y$ L-surface, draws the heatmap in Plotly, and adds the 2D
map to the PDF report. Simultaneous sharp components are also supported.
Performance improvements for very large grids and advanced sample-type-specific
validations are left for later phases.

\end{document}
